\chapter{数据假设与筛选：处理规则为什么需要验证}
\label{ch:reliability}

补值、删除、变换与筛选都在改变模型看到的信息。
本章把常见处理建议背后的假设拆开，避免把一个统计规则当作无条件的数据清洗标准。

\section{MCAR、MAR与MNAR：缺失依赖什么？}
\label{sec:reliability-mechanism}
考虑一个可能缺失的变量$X$、完整观测的变量$Z$，令$R=1$表示$X$被观测到。
下面用这个简化设置理解三类缺失机制，不把“随机”误读成没有任何规律。

\begin{center}
\begin{tabularx}{\textwidth}{P{29mm}P{64mm}Y}
\toprule
机制 & 条件关系 & 教学例子\\
\midrule
MCAR：完全随机缺失 & $P(R=1\mid X,Z)$为常数。 & 与记录内容无关的随机传输丢失。\\
MAR：随机缺失 & $P(R=1\mid X,Z)=P(R=1\mid Z)$。 & 保留率只由已记录的站点决定。\\
MNAR：非随机缺失 & 给定$Z$后，保留率仍依赖未观测的$X$。 & 收入本身越高，越不愿填写收入。\\
\bottomrule
\end{tabularx}
\end{center}

MAR允许缺失概率依赖已观测信息；它不是MCAR的同义词。
MNAR也不等于“缺失不是随机掷硬币”：关键是给定已观测信息后，是否仍依赖缺失部分。
一般多变量缺失模式比这里复杂，须针对可用观测和分析目标定义。

\begin{pitfall}[不能只看观测表就证明MAR]
缺失部分的真值通常不可见，不同机制可能产生相同的观测分布。
应结合采集流程、领域知识与敏感性分析；“没发现关系”不能证明缺失完全随机。
是否采用缺失指示变量、模型填补或专门缺失机制模型，仍需验证。
\end{pitfall}

\quick{缺失机制讨论数据怎样被观察；填补方法讨论如何处理已得到的表。两者不能凭名称一一配对。}

\clearpage
\section{完整记录分析：MAR下也可能改变总体}
\label{sec:reliability-deletion}
完整记录分析（Complete-case Analysis）只保留所需变量齐全的记录。
是否合适取决于估计目标、缺失机制与模型，不能说“缺失原因在其他特征中体现，就一定可以删”。

\begin{figure}[H]
\centering\begin{tikzpicture}[x=1cm,y=1cm]
\node[lab,left] at (0,.5) {总体100人};
\fill[Teal] (0,0) rectangle (5,1);\fill[Amber] (5,0) rectangle (10,1);
\node[white] at (2.5,.5) {A：50人，$X=20$};\node[white] at (7.5,.5) {B：50人，$X=80$};
\node[lab,left] at (0,-1.3) {完整记录50人};
\fill[Teal] (0,-1.8) rectangle (8,-.8);\fill[Amber] (8,-1.8) rectangle (10,-.8);
\node[white] at (4,-1.3) {A：40人};\node[white] at (9,-1.3) {B：10人};
\node[lab,below] at (5,-2.1) {每条横条各自归一化：保留率不同改变组别构成。};
\end{tikzpicture}

\caption{可见全部真值的教学构造。A组保留率0.8，B组保留率0.2，只依赖已观测组别$Z$。}
\end{figure}
\readfigure{总体两组各50人，A的$X=20$、B的$X=80$，总体均值为50。
按图中预期保留构成，只看完整记录的均值为$(40\times20+10\times80)/50=32$。
即使缺失满足本例MAR关系，不加调整的完整记录均值也偏离总体目标。}

如果知道目标总体两组各占一半，并且组内观测足以代表组内目标，
按目标组比例加权可恢复$0.5\times20+0.5\times80=50$。
现实中不能预先看到缺失真值，这种调整也依赖相应假设与足够覆盖。

MCAR下，完整记录可能仍是对总体的随机子样本，但会损失精度，
少量完整记录也不足以可靠训练复杂模型。某些条件模型在更弱机制下仍可有效估计，
因此也不能反过来说“只有MCAR才允许任何完整记录分析”。

\quick{先写清要估计的是总体均值、条件关系还是预测表现，再论删除与填补；对预测任务也要模拟实际缺失模式。}

\clearpage
\section{异常值：标记、核查，再决定处理}
\label{sec:reliability-outlier}
异常值（Outlier）是相对于某种分布、模型或规则显得少见的记录。
少见不等于错误；真实故障、罕见疾病或极端交易可能正是目标。

\begin{figure}[H]
\centering\begin{tikzpicture}[node distance=7mm]
\node[box] (a) {1};\node[box,right=of a] (b) {2};\node[box,right=of b] (c) {2};
\node[box,right=of c] (d) {3};\node[box,right=of d,draw=Amber,fill=Amber!8] (e) {100};
\node[box,right=16mm of e,text width=40mm] (summary) {均值：21.6\\中位数：2};
\draw[flow] (e)--(summary);
\node[lab,below=9mm of c,text width=68mm,align=center] {100很极端，但是否错误仍需查来源；\\统计标记不能替代业务核查。};
\end{tikzpicture}

\caption{五个数的例子；方框只列记录，不表示等距数轴。}
\end{figure}
\readfigure{均值被100明显拉动，中位数仍为2。若把100截到5，均值变为2.6；
这是修改了数据和估计对象，不能仅凭数字更“正常”就宣称处理正确。}

按本例线性分位数约定，$Q_1=2,Q_3=3$，$\mathrm{IQR}=1$。
常见箱线图围栏$[Q_1-1.5\mathrm{IQR},Q_3+1.5\mathrm{IQR}]=[0.5,4.5]$
会标记100；围栏是诊断启发式，不是测量错误的判决。
不同分位数定义、分组或时间范围可能改变标记结果。

\begin{center}
\begin{tabularx}{\textwidth}{P{29mm}Y}
\toprule
证据与选择 & 处理考虑\\
\midrule
确认录入或单位错误 & 按来源修正，无法修正时明确标记或排除理由。\\
真实但极端 & 保留目标事件；比较稳健损失、适当变换或分组模型。\\
暂时不确定 & 记录标记并做敏感性比较，不默认“标记最安全”或“取对数总有用”。\\
\bottomrule
\end{tabularx}
\end{center}

从数据估计的截断阈值与变换参数仍只在训练部分拟合；不要为了改善测试分数事后删除困难样本。

\clearpage
\section{方差筛选：变化大小不是预测价值}
\label{sec:reliability-variance}
方差阈值（Variance Threshold）是不使用标签的筛选方法。
近乎常数的特征在某些任务中贡献有限，但方差受到量纲影响，低方差也可能保留关键标签信息。

\begin{center}
\begin{tabular}{lrr}
\toprule
二元变量，取1比例$p$ & 分母为$n$的经验方差$p(1-p)$ & 含义\\
\midrule
$p=0.5$ & 0.25 & 0与1一样常见。\\
$p=0.1$ & 0.09 & 1较少见。\\
$p=0.01$ & 0.0099 & 1很少见，但可能重要。\\
\bottomrule
\end{tabular}
\end{center}

如果标签恰为$y=x$，即使$p=0.01$，该特征仍可完美区分正类。
仅因它低方差就删除，会丢掉全部检出正类的信息。
连续特征乘100后，方差乘10000，含义却未必改变；跨不同单位直接比较原方差不合理。

标准化为单位方差后，再用一个普通方差阈值筛选，也可能失去原来想识别的近常数差别。
应先明确筛选目标与处理顺序；二元指示变量、连续量与缺失指示不宜机械套同一阈值。

\begin{keyidea}[筛选也属于拟合]
阈值、特征列表、重要性排名或目标关联一旦从数据学习，
必须在交叉验证每折的训练部分重新确定，再将同一规则用于该折验证。
“这是预处理”不能成为在全部数据上先看标签筛选的理由。
\end{keyidea}

\clearpage
\section{卡方关联筛选：频数、期望与局限}
\label{sec:reliability-chi2}
对类别特征与类别目标，可用列联表比较观测频数与独立假设下的期望频数。
设总计$N$，则$E_{rc}=O_{r+}O_{+c}/N$，Pearson卡方统计量为
\[
 \chi^2=\sum_{r,c}\frac{(O_{rc}-E_{rc})^2}{E_{rc}}.
\]
这里求和的是表格单元，不是任意把连续取值当频数。

\begin{center}
\begin{tabular}{lrrr}
\toprule
40条记录 & $y=1$ & $y=0$ & 合计\\
\midrule
$x=1$ & 12 & 8 & 20\\
$x=0$ & 4 & 16 & 20\\
\midrule
合计 & 16 & 24 & 40\\
\bottomrule
\end{tabular}
\end{center}

本例期望表为$\left[\begin{smallmatrix}8&12\\8&12\end{smallmatrix}\right]$，
四格差的绝对值均为4，因此$\chi^2=20/3\approx6.667$。
使用无连续性校正的Pearson统计量；固定2×2类别设置下自由度为1。
对稀疏期望频数，渐近近似可能不好，应考虑合适的精确或置换方法。

统计关联不等于因果或独立测试上的预测收益。
单变量筛选可漏掉交互：XOR中每个独立均匀二元输入单独与标签无关联，两者一起却能决定标签。
大量特征反复检验还涉及多重选择，不能把筛选后的普通$p$值无条件解读。

\quick{某些软件的非负特征卡方筛选有专门数据约定；不可把已标准化的负数直接送入要求非负计数的接口。}

\clearpage
\section{向前选择：每一步都是一次模型比较}
\label{sec:reliability-forward}
向前逐步选择（Forward Selection）从空特征集或预设起点开始，
每次分别尝试加入一个尚未选择的特征，按约定的验证指标保留当前最好扩展。

\begin{center}
\begin{tabularx}{\textwidth}{P{29mm}Y}
\toprule
步骤 & 示例流程，四个候选特征\\
\midrule
第一轮 & 比较$\{x_1\},\{x_2\},\{x_3\},\{x_4\}$；假设本轮选$x_2$。\\
第二轮 & 比较$\{x_2,x_1\},\{x_2,x_3\},\{x_2,x_4\}$。\\
决定停止 & 按预定特征数或验证改进规则停止；不能使用最终测试分数控制。\\
\bottomrule
\end{tabularx}
\end{center}

取到两个特征需比较$4+3=7$个候选子集；空集基线另算。
若每组候选五折评估，至少涉及35次折内模型拟合，预处理与候选搜索也会增加开销。
在更多特征中，它比穷举所有子集便宜，但这是贪心搜索，不保证全局最佳子集。

若用同一批数据既选择特征又报告误差，结果会受到选择影响。
需要独立测试或合适的外层评估，把整个特征选择流程包在训练部分中。
强相关特征可能互换，交互也可能使单独加入时看不出收益。

\quick{特征筛选依赖模型和指标。某次没选上不代表该变量没有信息，也不代表它没有科学或因果意义。}

\clearpage
\section{Tomek链接与样本量经验的边界}
\label{sec:reliability-tomek}
Tomek链接是在给定距离下、互为最近邻且类别不同的一对样本。
这里假设没有距离并列；有并列时要说明选择约定。

\begin{figure}[H]
\centering\begin{tikzpicture}[x=1cm,y=1cm]
\node[box] (a) at (0,0) {$0:A$};\node[box] (b) at (2.6,0) {$1:A$};
\node[box,draw=Amber] (c) at (5.2,0) {$1.1:B$};\node[box,draw=Amber] (d) at (8.2,0) {$3:B$};
\node[box,draw=Amber] (e) at (11,0) {$3.2:B$};
\draw[flow,Quiet] (a)--(b);
\draw[{Stealth[length=2mm]}-{Stealth[length=2mm]},Teal,very thick] (b)--(c);
\draw[{Stealth[length=2mm]}-{Stealth[length=2mm]},Quiet,thick] (d)--(e);
\node[lab,above] at (3.9,.6) {异类，距离0.1：Tomek链接};
\node[lab,below] at (9.6,-.6) {同类，虽互近但不符合条件};
\end{tikzpicture}

\caption{五个一维数值及类别；箭头指最近邻，横向间距仅作示意，不按数值比例绘制。}
\end{figure}
\readfigure{1与1.1异类且互近，构成链接；3与3.2虽互近但同类，因此不构成链接。
B是本例多数类，若采用“只移除链接中多数类样本”规则，会移除1.1，A/B从2/3变为2/2。}

也有同时移除两端的清理约定，两者改变的数据不同。
清理可能去掉噪声或边界重叠，也可能丢失有效边界信息；不保证任何模型的测试表现提高。
仍只在训练折内使用，验证与测试应保持目标场景分布。

\begin{pitfall}[没有通用的每类5000条或总量1000万条门槛]
所需样本量取决于任务难度、噪声、分布、表示、模型、迁移与希望达到的误差。
不要把某本材料中的经验量级变成训练资格标准或人类水平的保证。
用同源边界正确的学习曲线、基线与目标误差检查数据是否仍限制效果。
\end{pitfall}

\noindent\textbf{参考：}Little与Rubin，\textit{Statistical Analysis with Missing Data}；
scikit-learn Feature Selection文档；imbalanced-learn Tomek Links文档。
图表与数值例子独立构造。

